\section{LINEAR REPRESENTATIONS AND GLOBAL LIE GROUPS.}

None of the work of which we have been speaking tackled frankly the problem of the definition and the study of global Lie groups. The first steps along this path go back to H. Weyl. He is inspired by two theories, which until then had developed independently: that of linear representations of semisimple complex Lie algebras, due to E. Cartan, and that of linear representations of finite groups, due to Frobenius and which had just been transposed to the orthogonal group by I. Schur using an idea of Hurwitz. This latter had shown {[}168{]} how invariants can be constructed for the orthogonal groups or the unitary groups by replacing the operation of the mean on a finite group by integration relative to an invariant measure. He had also remarked that in applying this method to the unitary group, invariants for the general linear group are obtained, a first example of the ``unitarian trick''. In 1924, I. Schur {[}279 e{]} used this procedure to show the complete reducibility of the representations of the orthogonal group O(n) and of the unitary group U(n), by the construction of a positive non-degenerate invariant hermitian form; he deduces from this, by the ``unitarian trick'', the complete reducibility of the holomorphic representations of O(n, C) and of SL(n, C), establishes orthogonality relations for the characters of O(n) and of U(n) and determines the characters of O(n). H. Weyl also extends this method to complex semisimple Lie algebras {[}331 b{]}. Given such an algebra g, he shows that it possesses a ``real compact form'' (which amounts to saying that it comes from an algebra \(g_{0}\) over R, whose adjoint group \(G_{0}\) is compact, by extending the scalars from R to C). Furthermore he shows that the fundamental group of \(G_{0}\) is finite, thus that the universal cover \(^{20}\) of \(G_{0}\) is compact. He deduces from this, by an appropriate adaptation of the procedure of Schur, the complete reducibility of the representations of g, and gives also, by global means, the determination of the characters of the representations of g. In a letter to I. Schur {[}331 a{]}, H. Weyl summarises the results of Cartan, that Schur did not know (cf.~{[}279 e{]}, p.~299, note at the foot of the page) and compares the two points of view: the method of Cartan supplies all the holomorphic representations of the simply connected group of the Lie algebra g; in the case of the orthogonal group, the representations of a cover in two sheets (later called the spinor group), which had escaped Schur's notice, are also obtained; from another side, Schur's method has the advantage of proving complete reducibility and of giving the characters explicitly.

After the work of H. Weyl, E. Cartan adopts an openly global point of view in his research on symmetric spaces and Lie groups. It is this point of view that is at the basis of his exposé of 1930 ({[}52 a{]}, v. I₂, pp.~1165-1225) of the theory of ``finite and continuous'' groups. There is found in particular the first proof of the global variant of the third fundamental theorem (existence of a Lie group with given Lie algebra); Cartan shows also that every closed subgroup of a real Lie group is a Lie group, which generalises a result of J. von Neumann on the closed subgroups of the linear group {[}324 b{]}. In this Memoir, von Neumann showed also that every continuous representation of a complex semisimple group is real analytic.

After this work, the theory of Lie groups in the ``classical'' sense (that is to say of finite dimension over R or C) is more or less fixed in its main lines. The first detailed exposé of it is given by Pontrjagin in his book on topological groups {[}253{]}; he keeps there to a point of view still fairly close to that of Lie, but in distinguishing carefully the local from the global. It is followed by the book of Chevalley {[}62 d{]} which contains also the first systematic discussion of the theory of analytic varieties and of the exterior differential calculus; the ``infinitesimal transformations'' of Lie appear there as fields of vectors and the Lie algebra of a Lie group G is identified with the space of fields of vectors invariant on the left under G. He leaves to one side the ``groupuscule'' aspect and the ``groups of transformations'' aspect.

